\begin{abstract}
A limited array can localize a dielectric inclusion without determining the surrounding material. We formulate the output as a set of possible centres and a set of readings that each surviving material class can produce under another illumination. Both are computed from the polarization currents of the nonlinear volume integral equation. Inertia bounds permit invisible current directions and local changes of the material interval admit signed low-rank updates. For two classes, a pair of nonnegative weights bounds the distance between their conditional response sets. Under strict feasibility and a coercivity condition the distance formula is exact for the two quadratic outer sets. A complementary construction gives a lower bound on the incident energy or number of simultaneously controlled transmitters. For specified outlines and contrast intervals in the discrete model, the data exclude 16,313 of 16,384 candidate centres across sixteen 32 by 32 specimens with independently varying exterior contrast. Two further specimens were examined on a 96 by 96 grid. Allowing the edge contrast of the annulus to vary leaves 973 of 9,216 centres possible. Across both grids a two-parameter template fit correctly located all eighteen centres and one fitted material lay inside the data ball. The set method found a compatible full-wave material at all eighteen centres. In two independent native-grid binary examinations, coherent illumination guarantees separation at energies for which every individual direction admits a common reading. At the chosen energy and deterministic error budget two simultaneous directions are necessary and sufficient in one examination. Coarsening the simulation grid changes the field by as much as 5.27 noise radii on the two native specimens. Measured microwave data from two laboratories show how receiver aperture and calibration limit the transfer of a beam chosen from measured fields.
\end{abstract}
